\chapter{}

Tutti i giorni Pietro l'attese: la rivedeva lì con le braccia sul tavolino. Ma la sensazione d'averla trovata soltanto e di non amarla cresceva. E non andava agli esami, quantunque ci pensasse continuamente e s'imaginasse, come in una allucinazione che lo spaventava, d'essere interrogato e di non rispondere.

Andò invece a cercare Ghìsola, con un'impazienza che lo faceva perfino piangere.

Ella stessa aprì l'uscio; e Pietro fu sorpreso di amare proprio lei nel momento che le chiese:

— Mi aspettavi?

Ella, per tenerlo a bada, rispose:

— Forse.

Allora, quantunque provasse una specie di contrarietà anche a parlare, gli venne detto:

— Non potremmo stare insieme nella strada? Sei sola?

Ghìsola riflettè; e poi rispose:

— Aspettami dinanzi alla Badia.

Pietro non ne provò nessun piacere, perchè il senso disagevole d’una menzogna indefinibile l’opprimeva. L’aspettò soltanto per non mancare a ciò che egli stesso le aveva chiesto.

Tirava vento; ma c’era dovunque il sole ardente di luglio. Per la strada di Bisarno, alcuni cipressi si movevano in fondo alla svolta. E pareva che la luce fosse continuamente cambiata dal vento. Olivi, in fila, sporgevano con i rami lungo un muro. E le loro chiome, d’un verde tenero, vi sbattevano sopra. E anche le loro ombre parevano chiome: a pena si distinguevano da quelle vere.

Ella venne a passi rapidi. Era senza cappello e portava al collo una catena con un cuoricino d’oro.

Pietro temette d’esser ridicolo dicendole che doveva tornare a Siena. Ma, infine, ella gli chiese, dopo aver camminato in silenzio, mentre egli le guardava sempre le mani:

— Quando vai via?

— Domani.

— Non ci vedremo più, dunque!

Egli, sorpreso di quella calma un poco scherzosa, chiese sospirando:

— Penserai sempre a me?

Allora Ghìsola rispose, con convinzione, quasi con ubbidienza:

— Sempre.

Poi lo guardò e vedendo la sua scontentezza, riprese:

— Tu pensi ch’io ti ami poco.

Egli, quantunque fosse vero, rispose:

— Mi fido di te.

Ghìsola, tenendo la testa bassa, risorrise; ma questa volta la bocca s’indugiò nell’atto piacevole.

Quella strada, dove il vento sollevava qua e là nembi di polvere bianca, senza farsi sentire, era così solitaria come non ci fosse mai passato nessuno. Ghìsola gli pareva bella in un altro modo, e più grassa. «Sì, anche così è vestita bene!» Ma egli non poteva levare gli occhi da quel cuoricino: glielo voleva portar via, perchè se no l’avrebbero guardata di più proprio nel petto.

Ghìsola se ne accorse e aspettava. Egli, allora, quando vide che se n’era accorta, le disse:

— Perchè lo tieni?

Ella arrossì e parve che volesse proteggere il cuoricino.

— L’hai comprato tu o ti è stato regalato?

— Regalato.

— Dimmi chi. Dimmelo subito.

Egli si soffermò dinanzi a lei, e l'obbligò a fare lo stesso.

— La mia sorella Lucia.

— Quanto tempo è?

— Anno, quando venne a trovarmi.

— E ti vuol bene?

— Lei sì, ma io no.

— Perchè?

— Non lo so....

— Perchè? Dimmelo. Se non lo dici a me!

— Non lo so. Non ci assomigliamo di carattere.

Egli pensò che potesse esser vero, perchè erano completamente dissimili anche di persona; e ne ebbe piacere. Ma nondimeno era geloso lo stesso anche della sorella. E le disse:

— Te ne comprerò uno io, e porterai il mio. Ossia il tuo, perchè niente è più mio. Sei contenta?

Ella aveva voglia da vero di ridere; ma, certo, non era il momento. Invece tornò indietro senza dir nulla. E siccome si mise a camminar lesta lesta come se avesse fatto tardi, egli chiese:

— T'aspetta quella donna?

— Sì, siamo stati imprudenti.

— Ma perchè dici così, se io ti amo da vero? Tu non devi preoccupartene.

Ella sorrise, ed allungò il passo senza rispondergli.

Pietro lasciò che arrivasse sola nella piazza; poi, facendo finta d'aspettare qualcuno, camminò lì d'intorno. Ma non c'era nessuno! Vide un cane che scappava con la groppa ossuta, ad arco.

Per la strada di Grassina, guardò la collinetta d'un verde pallido e sbiadito, tutta oliveti; con cipressi qua e là, mescolati, sottili.

Arrivato da quella svolta un tranvai, egli vi salì. Quando alzò gli occhi era già dentro Firenze, passato di poco la Barriera, sul Lungarno biancheggiante; e vide da quel punto tutti i campanili insieme.
